\documentclass[10pt,a4paper]{amsart}
\usepackage[margin=19mm]{geometry}
\usepackage{amsmath,amssymb,mathtools}
\usepackage[scr=rsfs,cal=euler]{mathalfa}
\usepackage{xfrac}
\usepackage[unicode,hidelinks,bookmarks=false]{hyperref}
\newcommand{\ie}{i.e.\ }
\newcommand{\cf}{cf.\ }
\newcommand{\resp}{respectively\ }
\newcommand{\Subsection}{\subsection}
\providecommand{\nobreakdash}{\nobreak\hbox{-}}
\setlength{\emergencystretch}{2em}
\multlinegap0pt
\usepackage{bm}
\usepackage{bbm}
\usepackage{dsfont}
\usepackage{xspace}
\usepackage{cancel}
\usepackage{mathtools}
\usepackage{amsthm}
\makeatletter
\@ifundefined{theorem}{\newtheorem{theorem}{Theorem}}{}
\@ifundefined{corollary}{\newtheorem{corollary}[theorem]{Corollary}}{}
\makeatother
\newcommand{\E}{\mathcal{E}}
\newcommand{\N}{\mathbb{N}}
\title[Multiplicative functions additive on partitions of $2k$ nonz]{Multiplicative functions additive on partitions of $2k$ nonzero squares}
\author{Jewel Mahajan}
\date{Source: arXiv:2606.29507v1 (CC BY 4.0)}
\begin{document}
\maketitle

\noindent\emph{Source and licence.} Multiplicative functions additive on partitions of $2k$ nonzero squares. Version arXiv:2606.29507v1 (CC BY 4.0), distributed under the Creative Commons Attribution 4.0 International licence, as recorded in the arXiv licence metadata for this exact version and shown on the abstract page https://arxiv.org/abs/2606.29507v1. Full rights notice: licenses/arxiv-2606.29507-CC-BY-4.0.txt.\par\smallskip

\noindent\textit{Editorial scope.} Counted unit: the Corollary whose source label is \texttt{cor:6sqand8sq} in the article (source0.tex lines 176--179, source label \texttt{cor:6sqand8sq}), delivered together with the article's own complete original proof of it (source0.tex lines 254--312, one \texttt{proof} environment of 59 lines). No mathematics is written, repaired or re-derived: statement and proof are the article's own text line for line. The proof is detached from the statement in the source: the article states the result at lines 176--179 and prints its proof at lines 254--312, and the proof environment's own header names the statement, so the association is the article's own. Required context retained uncounted: eq:main (lines 125--128); eq:reduced (lines 131--134); thm:kge5-intro (lines 150--152). Every definition, display and label of those lines is retained unchanged. Omitted from this package and not redistributed: 1--124, 129--130, 135--149, 153--175, 180--253, 313--327. Nothing inside a retained excerpt is omitted or altered: every retained body is delivered exactly as the article's lines are written, so no omission receipt of any kind accompanies this package. Citations: the retained excerpts carry no citation command at all, so no bibliography entry of the article is transcribed and no reference is redistributed. Reference adaptation: none was needed and none was made; every reference inside the delivered unit resolves inside it, so no unresolved reference or citation is produced. Printed slips kept literal: no printed word, symbol, hypothesis or number of the counted statement or of its proof was altered. Layout and class: the article's own class is \texttt{amsart} with packages including geometry, microtype, hyphenat, parskip, amsmath, amssymb, amsthm, mathtools, commath, relsize, color, graphicx, array, subfig; the wrapper is the lane's pinned \texttt{amsart} editorial wrapper at 10pt on A4, which restates the theorem-like environments the retained text needs and adds the article's own macro definitions that the retained text needs (\texttt{\textbackslash E}, \texttt{\textbackslash N}), with extra wrapper packages (bm, bbm, dsfont, xspace, cancel, mathtools). The delivered numbering is the unit's own. Assets: no retained span contains a figure environment, a graphics-inclusion command, a PDF-page inclusion, a table or a TikZ picture; no image file is copied into this package, the package declares no asset dependency and no image file is redistributed with it. Licence: version arXiv:2606.29507v1 of this article is distributed under the Creative Commons Attribution 4.0 International licence, as recorded in the arXiv licence metadata for this exact version and shown on the abstract page https://arxiv.org/abs/2606.29507v1. A full rights notice is delivered at licenses/arxiv-2606.29507-CC-BY-4.0.txt. Characters: every retained excerpt is pure ASCII, so the delivered engine, XeLaTeX with its default Latin Modern text font, reports no missing character.
\begin{equation}\label{eq:main}
   f\Bigl(\sum_{i=1}^{2k}x_i^2\Bigr)=\sum_{j=1}^{k}f\bigl(x_{2j-1}^2+x_{2j}^2\bigr),
   \qquad x_1,\dots,x_{2k}\ge 1 .
\end{equation}

\begin{equation}\label{eq:reduced}
   f\Bigl(\sum_{i=1}^{k}A_i\Bigr)=\sum_{i=1}^{k}f(A_i)
   \qquad\text{for all } A_1,\dots,A_k\in\E_2 ,
\end{equation}

\begin{theorem}\label{thm:kge5-intro}
Let $k \ge 5$ and let $f$ be multiplicative and satisfy~\eqref{eq:main}. Then either $f(n) = n$ for all positive integers $n$, or $f(n) = 0$ for all $n > 2k + 21$.
\end{theorem}

\begin{corollary}\label{cor:6sqand8sq}
Every integer $n \ge 20$ is a sum of $6$ nonzero squares, and every integer
$n \ge 22$ is a sum of $8$ nonzero squares.
\end{corollary}

\begin{proof}[Proof of Theorem~\ref{thm:kge5-intro}]
Since $f$ satisfies~\eqref{eq:main}, $f$ also satisfies~\eqref{eq:reduced}. Let $N_0(k) \coloneqq 2k + 21$. Fix $N > N_0(k)$ and set $M = N - (2k+2)$. Since $N > 2k+21$, we have $M \ge 20$. Therefore, by Corollary~\ref{cor:6sqand8sq}, $M$ is a sum of $6$ non-zero squares, that is,
$M = C_3 + C_4 + C_5$ with $C_3, C_4, C_5 \in \E_2$. Setting $C_6 = \cdots = C_k = 2 \in \E_2$
(we assign nothing when $k = 5$), we obtain
\[
    \sum_{i=3}^{k} C_i = M + 2(k-5) = (N - 2k - 2) + (2k - 10) = N - 12,
    \qquad C_3,\dots,C_k \in \E_2 .
\]
Using $12 = 2 + 10$ and $13 = 5 + 8$, we have
\[
    N = 2 + 10 + \sum_{i=3}^k C_i, \qquad N + 1 = 5 + 8 + \sum_{i=3}^k C_i .
\]
Since $2,5,8,10 \in \E_{2}$, applying~\eqref{eq:reduced} to both yields
\[
    f(N) = f(2) + f(10) + \sum_{i=3}^k f(C_i), \qquad
    f(N + 1) = f(5) + f(8) + \sum_{i=3}^k f(C_i).
\]
Let $c = f(5) + f(8) - f(2) - f(10)$. Therefore, for all $N > N_0(k)$, we have
\[
    f(N+1) - f(N) = c.
\]
Hence, for all $n \ge N_{0}(k)+2$, the telescoping sum gives
\begin{align}
    &f(n)-f(N_0(k) + 1)\\
    =&\sum_{N=N_{0}(k)+1}^{n-1}(f(N+1) - f(N)) =\sum_{N=N_{0}(k)+1}^{n-1} c=c(n-N_{0}(k)-1).
\end{align}
Setting $d = f(N_0(k) + 1) - c(N_0(k) + 1)$, we have $f(n) = cn + d$ for all $n \ge  N_0(k)+2$.


For distinct primes $p, q$ with $p > q \ge N_0(k)+2$, multiplicativity of $f$ gives
\[
    cpq + d = (cp + d)(cq + d) = c^2 pq + cd(p + q) + d^2.
\]
Rearranging the above equation, we have
\[
    (c - c^2)pq - cd(p + q) + (d - d^2) = 0.
\]
Fix a prime $q_0\ge N_0(k)+2$. Then, for all primes $p>q_0$, we have
\[
    ((c - c^2)q_{0} - cd)p + (d - d^2) -cd q_{0} = 0.
\]
Since the above equality holds for infinitely many primes $p$, both coefficients must vanish, yielding
\[
    (c - c^2)q_{0} - cd=0, \qquad -cd q_{0} + (d - d^2)=0.
\]
Since $q_0\ge N_0(k)+2$ was arbitrary, we have
\[
    (c - c^2)q - cd=0, \qquad -cd q + (d - d^2)=0,
\]
for infinitely many primes $q\ge N_0(k)+2$. These further yields
\[
    c = c^2, \qquad cd = 0, \qquad d = d^2,
\]
so $(c, d) \in \{(0,0), (0,1), (1,0)\}$.

When $(c,d) = (0,1)$, we have $f(n) = 1$ for all $n \ge N_0(k)+2$. We now choose $A_1, \dots, A_k \in \E_2$ such that each exceeds $N_0(k)+2$. Then~\eqref{eq:reduced} gives \[1=f(\sum_{i=1}^{k}A_i)=\sum_{i=1}^{k} f(A_i) = k \ge 5,\] a contradiction. Hence $(c,d) \in \{(0,0),(1,0)\}$, that is, either $f(N) = N$ for all $N > N_0(k)$, or $f(N) = 0$ for all $N > N_0(k)$.

We now prove that if $f(N) = N$ for all $N > N_0(k)$, then $f$ is indeed the identity function on $\N$. Let $n\in \N$ be arbitrary. Pick a prime $p$ such that $p>\max (n,N_{0}(k))$. In particular, $pn>N_{0}(k)$ and $(p,n)=1$. By multiplicativity, $pn=f(pn)=f(p)f(n)=pf(n)$, which yields $f(n)=n$. Since $n\in \N$ was arbitrary, this proves $f$ is the identity function on $\N$.
\end{proof}

\end{document}
